\section{Interpretation}\label{sec:results/interpretation}

\subsubsection*{Influence of Cluster Similarity on Generalization}
To assess how well the model generalizes to structurally dissimilar proteins, performance was evaluated across different levels of cluster similarity between training and validation data.\autoref{fig:cluster_similarity_performance} shows the model performance (\gls{auc} and Top-1 accuracy) as a function of cluster similarity between training and validation data. A clear relationship between similarity and generalization performance can be observed.

When the cluster similarity is high (i.e., proteins in the training and validation sets are very similar), the model achieves its best validation performance. This indicates that the model is able to effectively leverage learned representations when the validation data is drawn from a distribution closely aligned with the training data.

However, as the similarity constraint becomes stricter and the overlap between training and validation proteins decreases, performance drops noticeably. In particular, when enforcing low similarity (e.g., ensuring that proteins differ by at least $70\%$), both \gls{auc} and Top-1 scores degrade. This demonstrates a reduced generalization capability: the model struggles to transfer learned knowledge to structurally dissimilar proteins.

This behavior can be explained by the structure of the input space, as illustrated in \autoref{fig:sequenceEmbeddingHumology} and \autoref{fig:proteinVsPEGComparison}. The protein embeddings form well-separated clusters, indicating that proteins with similar sequences or properties are grouped closely together in the embedding space. Furthermore, these clusters correlate with experimental conditions, suggesting that the model can partially rely on protein identity as a proxy for predicting outcomes.

As a result, when training and validation data share similar clusters, the model benefits from this implicit clustering structure and performs well. In contrast, when validation proteins belong to distant or previously unseen clusters, the model cannot rely on these learned associations, leading to a drop in performance. This highlights that the model has learned cluster-specific patterns rather than fully generalizable representations, which explains the observed decrease in generalization power under low-similarity conditions.

\begin{figure}[h]
	\centering
	\begin{overpic}[width=\textwidth, tics=10]{../msc-thesis-code/data/dataExploration/plots/humology}
		\put(23,-2.5){\small (A)}
		\put(73,-2.5){\small (B)}
	\end{overpic}
	\vspace{0.005cm}
	\caption[Cluster Similarity]{The Figure shows \gls{auc} and Top 1 score in relation to cluster homology. The model has difficulties generalizing towards models that are completely different from the ones in the training dataset.}
	\label{fig:cluster_similarity_performance}
\end{figure} 

\subsubsection*{Ablation Study}
To better understand the contribution of individual input features to model performance, an ablation study was conducted. By systematically removing each feature, the relative importance of different information sources for predicting crystallization outcomes can be assessed.
\autoref{fig:ablationStudy} presents the results of the study. The evaluation is shown in terms of mean margin and \gls{auc} for both training and validation data. Note that this was done for all features. Interestingly, the binding sites vector proved to be an interference signal. This is likely due to the unfit model selection of the features. The binding sites have local context which means a model able to pick up local features such as a \gls{cnn} would have been more suitable. The model training was continued without this vector and performance increased. The following studies thus present not all features presented in \autoref{sec:methods/feature_engineering}. The complete study showing the inference signal is depicted in \autoref{fig:ablationStudyFull} in \autoref{app:results}.

Across both metrics, disabling the \textit{structure embedding} leads to the most significant drop in performance. In particular, both mean margin and \gls{auc} decrease sharply on the training as well as the validation set, indicating that structural information provides the most critical signal for distinguishing compatible from incompatible conditions. This suggests that the geometric and spatial features of the protein implicitly encoded in the embedding are essential for the models' capability to rank crystallization conditions.

Removing \textit{sequence information} also results in a noticeable performance degradation, especially compared to removing the sequence embedding alone. This indicates that explicit sequence-derived features contribute additional information beyond learned embeddings. In contrast, disabling the \textit{sequence embedding} leads to only a minor decrease, suggesting that its contribution is less critical in the presence of other features.

The removal of \textit{surface information} results in a moderate performance drop, particularly on the validation set, indicating that surface-level features provide useful but not dominant information. In this ablation study only one feature type was removed at a time. Meaning while surface features might be important they might also be incorporated in the sequence or structure embedding. Overall, the ranking of feature importance based on performance degradation is: structure embedding $>$ sequence information $>$ surface information $>$ sequence embedding. The low importance of sequence embedding suggests that the model is able to learn beyond the homology artifact explored earlier. 

In addition to the ablation study, a selection study was conducted to investigate the minimal feature set required for the model to maintain performance. The result can be is illustrated in \autoref{fig:dependencyStudy} in \autoref{app:results}. In contrast to the ablation setting this study isolates each feature group by training the model with only a single feature type at a time.

The results are consistent with the findings from the ablation study. In particular, structure embeddings emerge as the most informative feature. This reinforces their central role in the model.

Interestingly, sequence embeddings also retain a notable amount of predictive power when used alone, performing substantially better than several other individual feature groups. This observation contrasts with the ablation study, where removing sequence embeddings from the full feature set had only a minor impact on performance.

This discrepancy suggests an important interaction effect:
while the \gls{esm} embeddings do encode relevant information, much of this information appears to be redundant with other feature groups. Consequently, when all features are combined, the marginal contribution of sequence embeddings is limited. However, in isolation, they can still provide meaningful signal, indicating that they capture overlapping, but not entirely unique, information. The overall lower performance in the dependency study suggest that interactions between feature groups are important as well. 

Finally, the gap between training and validation performance persists across all ablations, but becomes most pronounced when important features are removed. This further highlights the role of these features in enabling robust generalization.

\begin{figure}[h]
	\centering
	\begin{overpic}[width=\textwidth, tics=10]{../msc-thesis-code/data/dataExploration/plots/ablationStudy_without_surface}
		\put(23,-2.5){\small (A)}
		\put(73,-2.5){\small (B)}
	\end{overpic}
	\vspace{0.005cm}
	\caption[Ablation Study]{Figure shows \gls{auc} and Top 1 score for the ablation study. Performance decreases strongest when removing structure features. Every removal (except sequence embedding) significantly (p < 0.01 ANOVA) reduces the performance metric.}
	\label{fig:ablationStudy}
\end{figure} 

\subsubsection*{Impact of Data Composition on Model Performance}\label{sec:results/interpretation/data_ratio}
To better understand how training set size and composition affect model performance, an additional ablation study was conducted in which only a fraction of the available training data was used for training. After constructing the training and validation split based on cluster similarity, the training set was further subsampled using proportions ranging from $10\%$ to $100\%$. A ratio of $0.1$, for example, indicates that only $10\%$ of the training proteins are retained.

Importantly, this subsampling was not performed uniformly at random. Instead, proteins were ranked according to their frequency in the dataset, and subsets were constructed either from the most frequently occurring proteins (\textit{top}) or from the least frequent proteins (\textit{bottom}). Because the dataset exhibits a strongly imbalanced protein frequency distribution, this distinction has a substantial effect on the resulting number of samples. As shown in the right panel of \autoref{fig:dataRatioStudy}, the top-ranked proteins already account for a large fraction of all samples at very small protein ratios, whereas the bottom-ranked proteins contribute only few samples until a large fraction of proteins is included. In other words, a small \textit{top} subset can contain approximately as many samples as a very large \textit{bottom} subset. The number of duplicated entries for the top 15 proteins which causes this effect can be seen in \autoref{fig:most_common_proteins} in \autoref{app:data}. 

The left and middle panels of \autoref{fig:dataRatioStudy} show the resulting model performance in terms of Top1 and \gls{auroc}. Performance appears to increase with higher protein ratio for both subset types. However, the two curves differ substantially: the \textit{top} subsets generally achieve stronger performance at small ratios than the \textit{bottom} subsets. This is expected, since the \textit{top} subset has the same variance in protein with a higher variance in condition per protein. However, the effect is not related to the number of samples.

To disentangle the effect of protein ratio from the effect of sample count, \autoref{fig:ratio_samples} in \autoref{app:results} visualizes the same performance metrics against the number of unique samples instead of the proportion of selected proteins. Under this representation, a pattern becomes visible: performance does not follow a simple monotonic relationship with the number of samples alone. In particular, the \textit{bottom} subsets can achieve comparable or even better performance than \textit{top} subsets at similar or smaller sample counts.

This suggests that the dominant factor is not merely the number of available samples, but rather the diversity of proteins represented in those samples. The \textit{top} subsets are dominated by a small number of highly frequent proteins and therefore contain many repeated observations of the same proteins under varying conditions. In contrast, the \textit{bottom} subsets contain a broader range of protein identities per sample budget, even though each individual protein is observed less often. The fact that these subsets remain competitive when performance is plotted against unique sample count indicates that exposure to a more diverse set of proteins is more beneficial than repeatedly observing the same proteins with variations.

\begin{figure}[h]
	\centering
	\begin{overpic}[width=\textwidth, tics=10]{../msc-thesis-code/data/dataExploration/plots/dataProportionPerformance}
		\put(18,-3){\small (A)}
		\put(50,-3){\small (B)}
		\put(83,-3){\small (C)}
	\end{overpic}
	\vspace{0.005cm}
	\caption[Data proportion Study]{Figure shows the relation of performance to data ratio for different selection modes. Proteins are sorted based on the number of times they appear in the \gls{pdb}. The mode "Top"/"Bottom" selects the most/least common proteins. While absolute count of unique numbers of samples is important (as seen by the difference between the modes) the variance of proteins seems to be more important. This is illustrated by the linear increase in performance different to the exponential one illustrated in the right plot.}
	\label{fig:dataRatioStudy}
\end{figure} 

Overall, the results demonstrate that model performance is governed more by protein-specific structure and representation than by the sheer number of samples, with repeated observations per protein being particularly valuable. At the same time, the reduced performance under low similarity highlights that the learned representations are not fully generalizable across diverse protein clusters, indicating a remaining dependency on protein-level patterns.